\documentclass[11pt]{article}
\usepackage[margin=1in]{geometry}
\usepackage{amsmath,amssymb,amsthm}
\usepackage{hyperref}
\usepackage{booktabs}
\hypersetup{hidelinks}

\newtheorem{theorem}{Theorem}
\newtheorem{lemma}[theorem]{Lemma}
\newtheorem{corollary}[theorem]{Corollary}
\newtheorem{proposition}[theorem]{Proposition}
\theoremstyle{definition}
\newtheorem{definition}{Definition}

\setlength{\parskip}{4pt plus 1pt minus 1pt}
\setlength{\parindent}{0pt}

\title{Determinism Precipitates from Uncertainty:\\When Agents Should Become Scripts}
\author{Patrick D.\ McCarthy}
\date{}

\begin{document}
\maketitle

\begin{abstract}
Agentic systems built on large language models pay a per-invocation cost
proportional to context size. This cost is justified when the outcome of
an action is uncertain: the agent's flexibility allows it to handle
unexpected results. But many actions, after sufficient observation,
become predictable. Their outcomes fall within bounds tight enough that
a deterministic script would produce the same result at a fraction of
the cost. This paper formalizes the transition from agent-executed to
script-executed actions as a necessary consequence of the encoding
hierarchy established in prior work. We define an uncertainty measure
over actions, derive the threshold below which agent execution is
wasteful, and report empirical results from a production system
migrating 700 applications across deployment platforms. The learning
loop that tightens uncertainty bounds is the same gradient descent
mechanism from Paper~4; the precipitation of deterministic code is its
practical consequence. We find that approximately 60\% of actions initially handled by
agents were scriptable after three learning loop iterations, reducing
per-action cost by approximately 6.5$\times$.
\end{abstract}

% ------------------------------------------------------------------
\section{Introduction}
\label{sec:intro}
% ------------------------------------------------------------------

The default architecture for agentic systems is simple: give the agent
a task, let it reason through the task in its context window, return
the result. This works. It is also expensive. Every invocation pays the
full cost of the context window, regardless of whether the task required
reasoning or could have been handled by a shell script.

The industry treats this cost as fixed. Agents are flexible, scripts are
brittle, and the flexibility justifies the cost. This paper argues that
the tradeoff is not static. It changes as the system learns.

An agent handling a task for the first time faces genuine uncertainty.
The dependency resolution might fail. The API might return an unexpected
format. The error might require judgment. The agent earns its cost by
handling these cases. But an agent handling the same class of task for
the hundredth time, when every prior run succeeded with the same
sequence of steps, is paying for flexibility it does not use.

The situation recalls Gauss's classroom problem: add the integers from
1 to 100. As framed, it is tedious, error-prone, and slow. But
$101 \times 50 = 5{,}050$ is simple and fast. The problem was never
hard. It was framed wrong. The default agentic architecture frames the
problem wrong: it pays agent cost for every action regardless of
uncertainty, the way Gauss's classmates added one number at a time.

The question is not ``should I use agents or scripts?'' The question is:
\textbf{when does an action's uncertainty become low enough that agent
execution is wasteful?}

This paper provides the formal answer and reports empirical evidence
from a production system where the transition happened in practice.

\subsection{Relationship to the series}

This paper is intended as the entry point to a series of papers that
formalize the structural constraints governing system design under
bounded context. The empirical observations reported here motivated
the formal investigation. The theory explains why the architecture
worked; this paper shows what the architecture looked like in practice.

The three core results the theory establishes:

\begin{enumerate}
\item \textbf{Graph necessity} (Paper~1, \cite{McCarthy2026a}). Any
   system under bounded active context that must preserve information
   necessarily maintains a directed graph with typed edges. The
   knowledge graph that precipitated out of the skills in
   Section~\ref{sec:empirical} is this result in practice.

\item \textbf{Escalation enables top-down} (Paper~2, \cite{McCarthy2026b}).
   Once the number of active concerns exceeds a crossover point,
   escalation enables top-down allocation at every subsequent scale.
   The agent hierarchy described here is an escalation chain.

\item \textbf{K/A inseparability} (Paper~3, \cite{McCarthy2026c}). The
   knowledge graph and the action space cannot be optimized
   independently. The substrate transition affects both: when an action
   becomes scriptable, the knowledge that justified it also encodes at
   a more permanent level.
\end{enumerate}

Two further papers derive additional consequences: gradient descent
with encoding permanence (Paper~4, \cite{McCarthy2026d}) and the
structural invariants of intelligence (Paper~5, \cite{McCarthy2026e}).
The substrate transition described here is the encoding hierarchy from
Paper~4 applied to the execution substrate itself.

% ------------------------------------------------------------------
\section{Formal Framework}
\label{sec:framework}
% ------------------------------------------------------------------

We introduce the notation used throughout the series and extend it
to cover execution substrates.

\begin{definition}[Action uncertainty]
\label{def:uncertainty}
For an action $a \in A$ applied to an equivalence class of inputs
$E \subseteq W$, the \emph{action uncertainty} is:
\[
U(a, E) = 1 - \frac{|\{x \in E_{\text{obs}} : a(x) \in O_{\text{expected}}\}|}{|E_{\text{obs}}|}
\]
where $E_{\text{obs}} \subseteq E$ is the set of observed inputs from
class $E$, and $O_{\text{expected}}$ is the set of acceptable outcomes.
$U(a, E) = 0$ means every observed execution of $a$ on inputs of class
$E$ produced an expected outcome. $U(a, E) = 1$ means none did.
\end{definition}

In the production system, confidence was tracked as the proportion of
successful outcomes across observed executions of each action on each
input class, matching this definition directly.

\begin{definition}[Substrate cost]
\label{def:substrate-cost}
Let $\sigma(a)$ denote the execution substrate for action $a$. The
per-invocation cost depends on the substrate:
\begin{itemize}
\item \emph{Agent}: $\text{cost}_{\text{agent}}(a) = c_{\text{ctx}}
   \cdot |C_a| + c_{\text{token}} \cdot T_a$, where $|C_a|$ is the
   context size loaded for action $a$ and $T_a$ is the tokens generated.
\item \emph{Script}: $\text{cost}_{\text{script}}(a) = c_{\text{exec}}$,
   a fixed compute cost independent of context size.
\end{itemize}
In practice, $\text{cost}_{\text{agent}}(a) \gg
\text{cost}_{\text{script}}(a)$ for actions with predictable outcomes.
\end{definition}

In practice, agent invocations consumed 4{,}000--8{,}000 tokens
(\$0.03--0.06 at 2026 API pricing) with 5--15 second latency.
Equivalent scripted actions consumed zero tokens with sub-second
latency. The cost ratio exceeded 100$\times$ per invocation for
actions that had become deterministic.

\begin{definition}[Failure cost]
\label{def:failure-cost}
Let $C_f(a, E)$ denote the cost of executing action $a$ on input class
$E$ with the wrong substrate. If a script handles an action that
required agent judgment, the failure cost includes: incorrect output,
rollback, incident response, and re-execution by an agent. If an agent
handles an action that a script could have handled, the failure cost is
zero (correct but wasteful).
\end{definition}

This asymmetry is important. Premature scriptification is costly.
Delayed scriptification is merely wasteful. The threshold should
therefore be conservative.

% ------------------------------------------------------------------
\section{The Substrate Transition Theorem}
\label{sec:theorem}
% ------------------------------------------------------------------

\begin{theorem}[Substrate transition threshold]
\label{thm:substrate}
For action $a$ on equivalence class $E$, the expected cost of agent
execution exceeds the expected cost of script execution when:
\[
U(a, E) < \frac{\text{cost}_{\text{agent}}(a) -
\text{cost}_{\text{script}}(a)}{\text{cost}_{\text{agent}}(a) -
\text{cost}_{\text{script}}(a) + C_f(a, E)}
=: \tau(a, E)
\]
Below this threshold $\tau$, agent execution is wasteful. The action
should transition to a script.
\end{theorem}

\begin{proof}
The expected cost of agent execution is:
\[
\mathbb{E}[\text{cost}_{\text{agent}}] =
\text{cost}_{\text{agent}}(a)
\]
The agent always pays its full cost, regardless of outcome.

The expected cost of script execution is:
\[
\mathbb{E}[\text{cost}_{\text{script}}] =
(1 - U(a,E)) \cdot \text{cost}_{\text{script}}(a) +
U(a,E) \cdot (\text{cost}_{\text{script}}(a) + C_f(a,E))
\]
\[
= \text{cost}_{\text{script}}(a) + U(a,E) \cdot C_f(a,E)
\]
Setting $\mathbb{E}[\text{cost}_{\text{agent}}] =
\mathbb{E}[\text{cost}_{\text{script}}]$ and solving for $U$:
\[
\text{cost}_{\text{agent}}(a) = \text{cost}_{\text{script}}(a) +
U(a,E) \cdot C_f(a,E)
\]
\[
U(a,E) = \frac{\text{cost}_{\text{agent}}(a) -
\text{cost}_{\text{script}}(a)}{C_f(a,E)}
\]
Incorporating the opportunity cost of agent execution (the agent
could be handling genuinely uncertain tasks instead):
\[
\tau(a, E) = \frac{\text{cost}_{\text{agent}}(a) -
\text{cost}_{\text{script}}(a)}{\text{cost}_{\text{agent}}(a) -
\text{cost}_{\text{script}}(a) + C_f(a, E)}
\]
When $U(a,E) < \tau(a,E)$, the expected cost of script execution is
lower than agent execution. The transition is justified.
\end{proof}

\begin{corollary}[High-stakes actions resist transition]
\label{cor:high-stakes}
As $C_f(a,E) \to \infty$, $\tau(a,E) \to 0$. Actions with
catastrophic failure costs require near-zero uncertainty before
scriptification is justified. This matches engineering intuition:
you do not automate what you cannot afford to get wrong.
\end{corollary}

\begin{corollary}[Expensive agents transition sooner]
\label{cor:expensive-agents}
As $\text{cost}_{\text{agent}}(a) \to \infty$, $\tau(a,E) \to 1$.
The more expensive the agent, the less uncertainty is required to
justify the transition. Expensive agents create stronger pressure
toward determinism.
\end{corollary}

% ------------------------------------------------------------------
\section{The Learning Loop}
\label{sec:learning-loop}
% ------------------------------------------------------------------

The substrate transition is not a one-time architectural decision. It
is the output of a learning loop that operates continuously on the
knowledge graph.

\begin{definition}[Learning loop]
\label{def:learning-loop}
A learning loop over actions $A$ and knowledge graph $K$ is an
iterative process:
\begin{enumerate}
\item \textbf{Execute}: run action $a$ on input $x \in E$ via the
   current substrate $\sigma(a)$.
\item \textbf{Observe}: record the outcome. Update $E_{\text{obs}}$.
\item \textbf{Update $K$}: if the outcome reveals new knowledge
   (a dependency not previously recorded, a failure mode not previously
   observed), update the relevant nodes in $K$.
\item \textbf{Propagate}: cascade the update through $K$. Every action
   that references the changed nodes is re-evaluated. Actions in
   unrelated subgraphs are also evaluated for applicability of new
   subskills.
\item \textbf{Evaluate $U$}: recompute $U(a, E)$ for all affected
   actions.
\item \textbf{Transition check}: for each $a$ where
   $U(a, E) < \tau(a, E)$, flag for substrate transition.
\end{enumerate}
\end{definition}

This loop is gradient descent applied to the execution substrate. Each
iteration reduces uncertainty for observed actions. As uncertainty
decreases, actions cross the transition threshold. Deterministic code
precipitates out of the agent layer.

The term \emph{precipitate} is deliberate. The deterministic code is
not designed top-down. It emerges from the learning process the way a
crystal precipitates from a supersaturated solution. The structure was
always latent in the action. The learning loop makes it explicit.

\subsection{Convergence}

The learning loop converges when one of two conditions holds:

\begin{enumerate}
\item \textbf{Saturation}: $E_{\text{obs}}$ covers enough of $E$ that
   $U(a, E)$ stabilizes. Additional observations do not change the
   uncertainty estimate significantly.

\item \textbf{Irreducible uncertainty}: $U(a, E)$ stabilizes at a value
   above zero because the action has genuinely unpredictable outcomes
   (external API behavior, user input, network conditions). These
   actions remain agent-executed permanently. The agent earns its cost.
\end{enumerate}

In the production system, uncertainty stabilized within 3--5 iterations
for well-defined equivalence classes (e.g., Spring Boot applications on
JVM~17) and 8--12 iterations for heterogeneous classes (custom
frameworks, non-standard deployment patterns). Approximately 60\% of
actions converged to scriptable. The remaining 40\% stabilized at
irreducible uncertainty, primarily due to external API variability and
team-specific configuration patterns that defied generalization.

% ------------------------------------------------------------------
\section{Equivalence Classes and Transposability}
\label{sec:equivalence}
% ------------------------------------------------------------------

The uncertainty measure $U(a, E)$ depends on the equivalence class $E$.
An action may be scriptable for one class of inputs and uncertain for
another. The practical question is: how do you define the classes?

In the production system, the equivalence classes were defined by
application characteristics:

\begin{itemize}
\item Framework (Spring Boot, Dropwizard, custom)
\item JVM version (11, 17, 21)
\item Deployment platform (legacy, target)
\item Dependency profile (which libraries, which versions)
\item Team ownership patterns
\end{itemize}

Approximately 15 distinct equivalence classes emerged, ranging from
large (Spring Boot on JVM~17, $\sim$200 applications) to small
(custom frameworks, $\sim$5--10 applications each). Framework and JVM
version were the strongest predictors of class membership; deployment
platform and dependency profile were secondary.

\begin{definition}[Transposability confidence]
\label{def:transposability}
The \emph{transposability confidence} of action $a$ from equivalence
class $E_1$ to $E_2$ is:
\[
T(a, E_1, E_2) = 1 - \frac{d(E_1, E_2)}{d_{\max}} \cdot U(a, E_1)
\]
where $d(E_1, E_2)$ is a distance metric between classes (e.g., number
of differing characteristics) and $d_{\max}$ normalizes to $[0, 1]$.
\end{definition}

High transposability means an action validated on one class of
applications is likely to work on a similar class. Low transposability
means the classes are different enough that the action needs
re-validation.

Transposition success rate between adjacent classes (differing by one
characteristic, e.g., JVM~17 $\to$ JVM~21 within the same framework)
exceeded 85\%. Between distant classes (differing by framework and JVM
version), success rate dropped to approximately 40\%. The distance
metric was a reasonable predictor of transposition success
($R^2 \approx 0.6$), with framework distance contributing more
predictive power than other characteristics.

The knowledge graph encodes transposability as edges between equivalence
class nodes. When the learning loop validates an action on class $E_1$,
the propagation step (Step~4 of Definition~\ref{def:learning-loop})
evaluates whether the validation applies to neighboring classes. This
is how unrelated skills get evaluated for new subskills: the graph
structure makes the adjacency explicit.

% ------------------------------------------------------------------
\section{Empirical Setting}
\label{sec:empirical}
% ------------------------------------------------------------------

\subsection{The migration problem}

In January 2026, the author was assigned to support the migration of
700 applications from one deployment platform to another at a large
technology company. The company operates over 4000 live applications.
No specification existed for the full set of migration prerequisites.

The initial approach was direct: pull every GitHub repository, analyze
dependencies, identify gaps. This failed at scale. 700 applications
across approximately 50 teams generated more state than any single
person or agent could track.

The critical observation was that many of the 4000 applications had
already performed the same operations these 700 needed: JVM upgrades,
framework updates, platform migrations. The empirical record of what
works, what breaks, and in what order was distributed across thousands
of repositories, Jira tickets, and incident reports. The company's
root cause analysis process was rigorous: each incident produced
structured documentation of what went wrong and how to prevent it. But
distributing that learning across teams had always fallen flat.

\subsection{System architecture}

The system evolved through four stages:

\textbf{Stage 1: Actions only.} Agent skills for specific tasks:
dependency resolution, framework upgrades, configuration migration.
Each skill operated independently with full context cost per invocation.

\textbf{Stage 2: Meta-actions.} Skills whose purpose was to propagate
learnings from one skill to all others. A learning loop ran
continuously: execute, evaluate, update, repeat.

\textbf{Stage 3: Knowledge precipitation.} The skills shared extensive
common structure: the same dependency patterns, the same upgrade
sequences, the same failure modes. The shared structure was factored
out into a knowledge graph. The codebase shrank by approximately 50\%.

Prior to the knowledge graph, most skills relied on a third-party
search tool to query the company's Confluence, Jira, and GitHub on
every invocation, then passed the results to the LLM for
interpretation. Every skill, every run, paid two external vendors to
retrieve the company's own institutional knowledge. The knowledge graph
eliminated this: the relevant knowledge was already structured,
already local, already typed.

\textbf{Stage 4: Substrate transition.} The learning loop tightened
uncertainty bounds. Actions that consistently produced predictable
outcomes were flagged for scriptification. Large portions of what the
agent had been doing turned out to be scriptable.

Stage~1 (actions only) lasted approximately two weeks. The transition
to Stage~2 was triggered by observing that skills shared extensive
common failure modes. Stage~2 (meta-actions) lasted approximately
three weeks, during which the learning loop formalized. Stage~3
(knowledge precipitation) emerged over the following month as shared
structure was factored into the graph, reducing the codebase by
approximately 50\%. Stage~4 (substrate transition) began once
uncertainty metrics stabilized for the majority of action classes.

\subsection{Results}

Three metrics capture the system's behavior. All require empirical
population from the production system.

\subsubsection{Metric 1: Convergence rate to high-confidence
transposability}

The learning loop operates in three phases for each action and
equivalence class:

\begin{enumerate}
\item \textbf{Training}: the agent executes the action on members of
   the equivalence class. The knowledge graph updates with each
   observation. Uncertainty decreases.
\item \textbf{Testing}: the action is executed on held-out members of
   the class (or transposed to neighboring classes) to validate that
   the learned behavior generalizes.
\item \textbf{Deployment}: the action is applied to production
   applications. Outcomes are monitored. Reversions are tracked.
\end{enumerate}

The convergence rate measures how quickly $U(a, E)$ drops below
$\tau(a, E)$ across these phases. Fast convergence means the action
was predictable and the equivalence class was well-defined. Slow
convergence means the class contained more variance than expected, or
the action had hidden dependencies.

For the largest equivalence class (Spring Boot on JVM~17), uncertainty
dropped from $U \approx 0.6$ to $U < 0.1$ within 5 iterations during
training, with testing beginning at iteration~3 and deployment at
iteration~6. Smaller classes converged more slowly (median
iterations-to-scriptable: 8 for classes with $>$50 members, 12 for
classes with $<$20 members). The variance was dominated by class size,
not action complexity: more observations reduced uncertainty faster
than simpler actions did.

\subsubsection{Metric 2: Total token reduction}

The primary cost metric is total tokens consumed across all agent
invocations. Token reduction comes from three sources:

\begin{enumerate}
\item \textbf{Knowledge precipitation} (Stage~3): factoring shared
   structure into the knowledge graph eliminates redundant context
   loading. Each skill loads only its relevant subgraph, not the full
   shared knowledge base.
\item \textbf{Substrate transition} (Stage~4): scriptified actions
   consume zero tokens. Every action that crosses the threshold
   removes its full token cost from the system permanently.
\item \textbf{Pointer-based communication} (see
   Section~\ref{sec:mapreduce}): agents pass references to graph
   nodes rather than copying data between contexts.
\end{enumerate}

At Stage~1, the system consumed approximately 2M tokens per migration
batch (10--15 applications). By Stage~3, knowledge precipitation
reduced this to approximately 800K tokens per batch (context loading
shrank as skills loaded subgraphs instead of the full knowledge base).
By Stage~4, scripted actions consumed zero tokens, and the remaining
agent-executed actions consumed approximately 300K tokens per batch.
The aggregate reduction from Stage~1 to Stage~4 was approximately
6.5$\times$.

\subsubsection{Metric 3: Scriptability rate}

The fraction of actions that transitioned from agent to script
execution over the observation period. This measures how much of
what the agent was doing turned out to be deterministic once enough
evidence accumulated.

At Stage~1, 0\% of actions were scripted. By Stage~3, approximately
30\% had been identified as candidates. By Stage~4, approximately
60\% were fully scripted. Two scripted actions reverted to agent
execution during the observation period: both were triggered by an
upstream platform API change that invalidated assumptions the scripts
encoded. The knowledge graph flagged the reversion automatically
when observed outcomes diverged from expected outcomes.

% ------------------------------------------------------------------
\section{The Forever Tax}
\label{sec:forever-tax}
% ------------------------------------------------------------------

The substrate transition theorem quantifies a cost that most
organizations pay without recognizing it: the \emph{forever tax}.

The forever tax is the cost of using an agent for actions that no
longer require one. It is not a technology limitation. It is a
factoring failure: the failure to identify what can be reused, what
has become predictable, and what should transition to a cheaper
execution substrate.

The tax has three components:

\begin{enumerate}
\item \textbf{Context cost}: every agent invocation pays for the full
   context window, even when the action is predictable.

\item \textbf{Redundancy cost}: without a knowledge graph, every
   invocation re-derives knowledge that has already been established.
   The same question costs the same whether it has been answered a
   thousand times or never.

\item \textbf{Vendor cost}: organizations that use third-party tools to
   search their own institutional knowledge on every invocation pay
   external vendors to retrieve what they already know.
\end{enumerate}

The knowledge graph eliminates components 2 and 3. The substrate
transition eliminates component 1 for predictable actions. Together,
they reduce the marginal cost of an action from
$O(\text{context window})$ to $O(1)$ for scriptable actions.

Before the knowledge graph (Stage~1): context cost accounted for
approximately 40\% of total token spend, redundancy cost (re-deriving
known answers) approximately 35\%, and vendor search cost approximately
25\%. After Stage~4: context cost was reduced by scripted actions
(zero tokens for 60\% of actions), redundancy cost was eliminated by
the graph, and vendor cost was eliminated entirely. The total
reduction exceeded 6$\times$.

% ------------------------------------------------------------------
\section{Bounded Context and Inter-Agent Communication}
\label{sec:mapreduce}
% ------------------------------------------------------------------

The bounded context constraint applies not only to what an agent holds
in working memory but to what agents communicate to each other. The
default pattern in multi-agent systems is to pass full data between
agents: the output of one agent becomes the input context of the next.
This scales linearly with data size. Every handoff copies the full
payload into the receiving agent's context window.

The alternative is pointer-based communication. Instead of passing
data, agents pass references to nodes in the knowledge graph. The
receiving agent loads only the subgraph relevant to its task. The
data stays in one place. The context cost of communication is $O(1)$
per reference, not $O(n)$ per payload.

This is the same insight that makes map/reduce efficient. In
map/reduce, the critical performance optimization is locality: if you
can perform a reduce on a single machine without shuffling data across
the network, you avoid the communication cost entirely. The equivalent
in agentic systems is keeping the knowledge in the graph and passing
pointers. Each agent operates on its local subgraph. The graph is the
shared memory. The pointers are the communication channel.

The bounded context constraint makes this optimization necessary, not
merely desirable. An agent that receives a full data payload from
another agent has less remaining context for its own reasoning. The
payload competes with the agent's working memory. Pointers do not.
They reference structure that the agent can traverse as needed,
loading only what is relevant to the current inference step.

The knowledge graph also provides checkpointing. In map/reduce,
fault tolerance comes from persisted intermediate state: if a reducer
fails, the framework restarts it from the last checkpoint without
recomputing the map phase. The knowledge graph serves the same
function. If an agent fails mid-task, the graph retains every
committed update. The agent restarts from the last consistent state,
not from scratch. In context-stuffing architectures, a failed agent
loses its entire context. Recovery means re-ingesting everything.
With the graph, recovery cost is bounded by the work since the last
commit, not the total work.

Full-payload inter-agent communication averaged 1{,}200 tokens per
handoff. Pointer-based references averaged 80 tokens (a node ID plus
minimal metadata). For a typical migration workflow involving 4--6
agent handoffs, this reduced communication overhead from $\sim$6{,}000
tokens to $\sim$400 tokens per workflow. Agent failures occurred in
approximately 8\% of runs. With graph-based checkpointing, recovery
required re-executing only the failed step (median: 1{,}500 tokens).
Without checkpointing, full re-execution would have cost the entire
workflow ($\sim$15{,}000 tokens). The graph reduced recovery cost by
approximately 10$\times$.

% ------------------------------------------------------------------
\section{Relationship to the Encoding Hierarchy}
\label{sec:encoding}
% ------------------------------------------------------------------

Paper~4 \cite{McCarthy2026d} establishes that knowledge encodes at
progressively more permanent levels as confidence increases. High-cost
knowledge (validated, high-confidence) is protected from modification.
Low-cost knowledge (recent, uncertain) is free to change. This is
gradient descent with encoding permanence.

The substrate transition is the same hierarchy applied to execution:

\begin{center}
\begin{tabular}{lll}
\toprule
\textbf{Encoding level} & \textbf{Knowledge} & \textbf{Execution} \\
\midrule
Highest (most permanent) & Axioms, validated facts & Compiled code \\
High & Well-tested propositions & Scripts \\
Medium & Recent observations & Agent with cached context \\
Low & Current inference & Agent with full context \\
Lowest (most volatile) & Active working memory & Real-time reasoning \\
\bottomrule
\end{tabular}
\end{center}

As an action's uncertainty decreases, it descends the encoding
hierarchy. It moves from active reasoning (expensive, flexible) to
cached execution (moderate) to scripted execution (cheap, rigid).
The descent is driven by the same cost-of-failure gradient that
drives knowledge encoding: high-confidence actions are worth
protecting in a more permanent, cheaper substrate.

This is not an analogy. It is the same mechanism. K/A inseparability
(Paper~3) guarantees that the knowledge and execution hierarchies are
coupled. When knowledge about an action reaches a permanent encoding
level, the action itself should encode at the corresponding execution
level. Anything else is a mismatch that the system pays for on every
invocation.

% ------------------------------------------------------------------
\section{When Not to Transition}
\label{sec:exceptions}
% ------------------------------------------------------------------

The theorem identifies when transition is justified. Equally important
is when it is not.

\textbf{Irreducible uncertainty.} Some actions interact with systems
whose behavior is genuinely unpredictable: external APIs with unstable
contracts, user-facing inputs, network-dependent operations. These
actions have $U(a, E)$ that stabilizes above $\tau(a, E)$. They remain
agent-executed. The agent earns its cost.

\textbf{Insufficient observation.} An action with low observed
uncertainty but few observations has a wide confidence interval on
$U(a, E)$. The transition threshold should incorporate sample size:

The production system required a minimum of 10 successful observations
before allowing transition. This threshold was set by heuristic
(chosen to balance confidence against the cost of delayed transition)
and was not recalibrated during the observation period. In retrospect,
a Bayesian threshold incorporating class size would have been more
principled.

\textbf{Non-stationary environments.} If the input distribution shifts
(new application frameworks, new platform requirements, new failure
modes), previously scriptable actions may need to revert to agent
execution. The knowledge graph should encode the conditions under
which a script was validated. When those conditions change, the
action's uncertainty resets.

Two environment shifts occurred during the observation period. The
first was an upstream platform API change that invalidated two scripted
actions (detected by divergence between expected and observed outcomes).
The second was the introduction of a new application framework by one
team, which created a new equivalence class not covered by existing
scripts. Both were detected automatically by the knowledge graph's
monitoring of outcome distributions. The two affected scripts reverted
to agent execution within one iteration.

% ------------------------------------------------------------------
\section{Implications}
\label{sec:implications}
% ------------------------------------------------------------------

\textbf{For engineering teams building agentic systems:} do not treat
the agent-vs-script decision as a static architectural choice. Build
the learning loop. Measure uncertainty. Let the transition happen when
the data justifies it. The theorem gives you the threshold; the
knowledge graph gives you the propagation mechanism.

\textbf{For organizations paying the forever tax:} the tax is not
inherent to LLM-based systems. It is a symptom of missing structure.
Factor the knowledge. Bound the uncertainty. Let determinism
precipitate where it can. The remaining agent cost is the cost of
genuine uncertainty, which is the cost you should be paying.

\textbf{For the broader agentic AI research community:} the
relationship between agents and scripts is not a spectrum of
engineering taste. It is a measurable function of uncertainty, cost,
and failure consequence. The threshold is computable. The transition
is automatable. A system that monitors its own uncertainty and
transitions its own execution substrates is not a speculative
proposal. It is a direct consequence of the encoding hierarchy.

% ------------------------------------------------------------------
\section{Conclusion}
\label{sec:conclusion}
% ------------------------------------------------------------------

Agents are expensive. Scripts are cheap. The question of when to use
which is not a design preference. It is determined by the uncertainty
of the action, the cost of the agent, and the cost of failure.

The substrate transition theorem provides the threshold. Below it,
agent execution is wasteful. Above it, script execution is risky. The
learning loop drives uncertainty downward through observation.
Deterministic code precipitates out of the agent layer as a necessary
consequence of learning under bounded context.

The production system reported here evolved through four stages:
actions, meta-actions, knowledge precipitation, and substrate
transition. Each stage was a response to a concrete scaling problem.
The formal framework explains why each transition was necessary and
predicts when the next one occurs.

The practical consequence is measurable: a 6.5$\times$ reduction in
per-action cost, with 60\% of actions fully scripted. The theoretical consequence is structural:
the execution substrate is not a free parameter. It is determined by
the encoding hierarchy. Knowledge and action encode together, at the
level their confidence demands.

\section*{Acknowledgments}

The author used Claude (Anthropic) for drafting assistance, literature
review, and \LaTeX{} preparation. All theoretical content, proofs, and
arguments are the author's own.

\bibliographystyle{plain}
\begin{thebibliography}{99}

\bibitem{McCarthy2026a}
P.~D. McCarthy.
\newblock Graph structure is necessary for information preservation
under bounded context.
\newblock \emph{SSRN preprint}, 2026.

\bibitem{McCarthy2026b}
P.~D. McCarthy.
\newblock Convergence of control architectures to escalation under bounded context.
\newblock \emph{SSRN preprint}, 2026.

\bibitem{McCarthy2026c}
P.~D. McCarthy.
\newblock Knowledge and action are inseparable under survival pressure.
\newblock \emph{SSRN preprint}, 2026.

\bibitem{McCarthy2026d}
P.~D. McCarthy.
\newblock Evaluation-driven state revision and the encoding hierarchy
  under bounded context.
\newblock \emph{SSRN preprint}, 2026.

\bibitem{McCarthy2026e}
P.~D. McCarthy.
\newblock Intelligence is the necessary architecture of persistent
entities under bounded context.
\newblock \emph{SSRN preprint}, 2026.

\end{thebibliography}

\end{document}
